\documentclass{article}
\usepackage{amsmath, amssymb, listings, xcolor}
\usepackage{hyperref}

\usepackage{color}
\usepackage{listings}
\definecolor{GrayCodeBlock}{RGB}{241,241,241}
\definecolor{BlackText}{RGB}{110,107,94}
\definecolor{RedTypename}{RGB}{182,86,17}
\definecolor{GreenString}{RGB}{96,172,57}
\definecolor{PurpleKeyword}{RGB}{184,84,212}
\definecolor{GrayComment}{RGB}{170,170,170}
\definecolor{GoldDocumentation}{RGB}{180,165,45}
\lstdefinelanguage{rust}
{
    columns=fullflexible,
    keepspaces=true,
    breaklines=true,
    frame=single,
    framesep=0pt,
    framerule=0pt,
    framexleftmargin=4pt,
    framexrightmargin=4pt,
    framextopmargin=5pt,
    framexbottommargin=3pt,
    xleftmargin=4pt,
    xrightmargin=4pt,
    backgroundcolor=\color{GrayCodeBlock},
    basicstyle=\ttfamily\color{BlackText},
    keywords={
        true,false,
        unsafe,async,await,move,
        use,pub,crate,super,self,mod,
        struct,enum,fn,const,static,let,mut,ref,type,impl,dyn,trait,where,as,
        break,continue,if,else,while,for,loop,match,return,yield,in
    },
    keywordstyle=\color{PurpleKeyword},
    ndkeywords={
        bool,u8,u16,u32,u64,u128,i8,i16,i32,i64,i128,char,str,
        Self,Option,Some,None,Result,Ok,Err,String,Box,Vec,Rc,Arc,Cell,RefCell,HashMap,BTreeMap,
        macro_rules
    },
    ndkeywordstyle=\color{RedTypename},
    comment=[l][\color{GrayComment}\slshape]{//},
    morecomment=[s][\color{GrayComment}\slshape]{/*}{*/},
    morecomment=[l][\color{GoldDocumentation}\slshape]{///},
    morecomment=[s][\color{GoldDocumentation}\slshape]{/*!}{*/},
    morecomment=[l][\color{GoldDocumentation}\slshape]{//!},
    morecomment=[s][\color{RedTypename}]{\#![}{]},
    morecomment=[s][\color{RedTypename}]{\#[}{]},
    stringstyle=\color{GreenString},
    string=[b]"
}

\begin{document}


\section{Semantics of \texttt{verify!}}

The \texttt{verify!} macro is responsible for checking whether a given proof correctly satisfies the circuit constraints

\begin{lstlisting}[language=rust, caption={Macro Definition}]
verify!(circuit_expr, [public_inputs...], proof_expr)
\end{lstlisting}

The macro expands into:
\begin{lstlisting}[language=rust]
Verifier::verify(&circuit_expr, &[public_inputs...], &proof_expr)
\end{lstlisting}



\begin{lstlisting}[language=rust, caption={Verification Function}]
pub fn verify(circuit: &Circuit, public_inputs: &[FieldElement], proof: &Proof) -> bool {
    if !Self::check_commitments(&proof.commitments, &proof.revealed_witness) {
        return false;
    }
    let mut value_map: HashMap<String, FieldElement> = HashMap::new();
    for (name, (val, _blinding)) in &proof.revealed_witness {
        value_map.insert(name.clone(), val.clone());
    }
    for (i, sig) in circuit.inputs.iter().enumerate() {
        if let Some(pub_val) = public_inputs.get(i) {
            if let Signal::Input(ref name) = sig {
                value_map.insert(name.clone(), pub_val.clone());
            }
        }
    }
    if !Self::constraints_satisfied(&circuit.constraints, &mut value_map) {
        return false;
    }
    true
}
\end{lstlisting}




\subsection{Background}

\begin{lstlisting}[language=rust]
pub struct Circuit {
    pub inputs: Vec<Signal>,
    pub outputs: Vec<Signal>,
    pub constraints: Vec<Constraint>,
}
\end{lstlisting}

\begin{lstlisting}[language=rust]
pub struct Proof {
    pub commitments: HashMap<String, Commitment>,
    pub revealed_witness: HashMap<String, (FieldElement, FieldElement)>,
}
\end{lstlisting}

\begin{equation}
\mathcal{V} : C \times I \times P \to \{ \text{true}, \text{false} \}
\end{equation}
where:
\[
\begin{aligned}
    &\mathbb{C} \text{ is a circuit.} \\
    &\mathbb{I} \text{ is a set of public inputs.} \\
    &\mathbb{P} \text{ is a proof (commitments + revealed witness values).} \\
    &\text{The function outputs } \texttt{true} \text{ if verification succeeds, and } \texttt{false} \text{ otherwise.}
\end{aligned}
\]

\section{Semantics}

\subsection{Extract Circuit Components}
Given a circuit $C$, extract:
\begin{align*}
    \mathbb{I} & : \text{set of public input signals},\\
    \mathbb{O} & : \text{set of output signals},\\
    \mathbb{C} & : \text{set of constraints}.
\end{align*}

Semantics:
\begin{equation}
    C = (\mathbb{I}, \mathbb{O}, \mathbb{C})
\end{equation}

\subsection{Extract Proof Components}
Given a proof $P$, extract:
\begin{align*}
    \mathbb{K} & : \text{set of commitments},\\
    \mathbb{R} & : \text{set of revealed witness values}.
\end{align*}

Semantics:
\begin{equation}
    P = (\mathbb{K}, \mathbb{R})
\end{equation}

\subsection{Check Commitments}

\begin{lstlisting}[language=rust, caption={Check Commitments}]
fn check_commitments(
    commitments: &HashMap<String, Commitment>,
    revealed_witness: &HashMap<String, (FieldElement, FieldElement)>,
) -> bool {
    if commitments.len() != revealed_witness.len() {
        return false;
    }
    for (name, (val, blind)) in revealed_witness.iter() {
        let recomputed = Commitment::new(val, blind);
        if commitments.get(name).map_or(true, |c| recomputed.committed_hash != c.committed_hash) {
            return false;
        }
    }
    true
}
....
impl Commitment {
    /// C = SHA256( val_bytes || blind_bytes || prime_bytes )
    pub fn new(val: &FieldElement, blinding: &FieldElement) -> Self {
        let val_bytes = val.value.to_le_bytes();
        let blind_bytes = blinding.value.to_le_bytes();

        let mut hasher = Sha256::new();
        hasher.update(&val_bytes);
        hasher.update(&blind_bytes);
        hasher.update(&val.prime.to_le_bytes());

        let result = hasher.finalize();
        Self {
            committed_hash: result.to_vec(),
        }
    }
}

\end{lstlisting}

For each revealed witness $(s, w_s, r_s) \in \mathbb{R}$, verify:
\begin{equation}
    H(w_s, r_s, p) = \mathbb{K}(s)
\end{equation}
where $H$ is the cryptographic hash function.
\subsection{Construct Value Map}

Build a value map \(\sigma\) mapping signal names to field elements.

\begin{itemize}
    \item \textbf{Add public inputs:}
    \[
        \sigma(s) = v, \quad \forall s \in \mathbb{I}, v \in \mathbb{I}
    \]
    \item \textbf{Add revealed witness values:}
    \[
        \sigma(s) = w_s, \quad \forall (s, w_s, r_s) \in \mathbb{R}
    \]
\end{itemize}

This results in a mapping:

\[
\sigma : S \to \mathbb{F}_p
\]

\subsection{Evaluate Constraints}

\begin{lstlisting}[language=rust, caption={Evaluate Constraints}]
fn constraints_satisfied(
    constraints: &[Constraint],
    value_map: &mut HashMap<String, FieldElement>,
) -> bool {
    for constraint in constraints {
        let left_val = Self::eval_signal(&constraint.left, value_map).unwrap();
        let right_val = Self::eval_signal(&constraint.right, value_map).unwrap();
        let result = match constraint.operation {
            Operation::Add => left_val.add(&right_val),
            Operation::Mul => left_val.mul(&right_val),
            Operation::Sub => left_val.sub(&right_val),
            Operation::Eq => {
                if !left_val.equals(&right_val) {
                    return false;
                }
                left_val.clone()
            }
        };
        match &constraint.output {
            Signal::Output(ref nm) | Signal::Witness(ref nm) => {
                value_map.insert(nm.clone(), result);
            }
            _ => {}
        }
    }
    true
}
\end{lstlisting}

For each constraint $(s_1 \oplus s_2 \to s_3) \in \mathbb{C}$:
\begin{align*}
v_1 &= \sigma(s_1), \quad v_2 = \sigma(s_2),\\
v_3 &= v_1 \oplus v_2 \mod p,\\
    \sigma(s_3) &\stackrel{?}{=} v_3.
\end{align*}
If any constraint fails, return \texttt{false}.

\subsection{Return Verification Result}
\begin{equation}
\mathcal{V}(C, I, P) = \begin{cases} 
    \text{true}, & \text{if commitments and constraints hold} \\ 
    \text{false}, & \text{otherwise} 
\end{cases}
\end{equation}


\section{Using Inference Rules}

We define \texttt{verify!} as a judgment:

\[
C, I, P \Downarrow \text{true} \quad \text{(if verification succeeds)}
\]

\[
C, I, P \Downarrow \text{false} \quad \text{(if verification fails)}
\]

Where:

\begin{itemize}
    \item $\mathbb{C}$ is the circuit (constraints, inputs, outputs).
    \item $\mathbb{I}$ is the set of public inputs.
    \item $\mathbb{P}$ is the proof (commitments + revealed witnesses).
    \item The output is \texttt{true} if verification succeeds, otherwise \texttt{false}.
\end{itemize}

\subsection{(1) Commitment Check}

\[
\frac{ \forall (s, w_s, r_s) \in P.\mathbb{R}, \quad H(w_s, r_s, p) = P.\mathbb{K}(s) }
{ P \Downarrow_{\text{commit}} \text{true} }
\]

If all revealed witness values $(w_s, r_s)$ correctly match their commitments $H(w_s, r_s, p)$, the commitment check passes.

Otherwise:

\[
\frac{ \exists (s, w_s, r_s) \in P.\mathbb{R}, \quad H(w_s, r_s, p) \neq P.\mathbb{K}(s) }
{ P \Downarrow_{\text{commit}} \text{false} }
\]

If any commitment check fails, verification fails immediately.


\subsection{(2) Construct Value Map $\sigma$}

\[
\frac{ \sigma = \{ s \mapsto v \mid s \in \mathbb{I}, v \in I \} \cup \{ s \mapsto w_s \mid (s, w_s, r_s) \in P.\mathbb{R} \} }
{ (C, P) \Downarrow_{\text{map}} \sigma }
\]

The value map $\sigma$ is built by:

\begin{itemize}
    \item Assigning public input values $v$ to input signals $s$.
    \item Assigning revealed witness values $w_s$ to witness signals $s$.
\end{itemize}

This map $\sigma$ will be used for constraint evaluation.


\subsection{(3) Constraint Satisfaction}

For each constraint $(s_1 \oplus s_2 \to s_3) \in \mathbb{C}$:

\[
\frac{ \sigma(s_1) = v_1, \quad \sigma(s_2) = v_2, \quad v_3 = v_1 \oplus v_2 \mod p, \quad \sigma(s_3) = v_3 }
{ \mathbb{C}, \sigma \Downarrow_{\text{constraints}} \text{true} }
\]

Meaning:

\begin{enumerate}
    \item Values of $s_1$ and $s_2$ are retrieved from $\sigma$.
    \item The operation $\oplus$ is applied to compute $v_3$.
    \item The computed result $v_3$ matches $\sigma(s_3)$.
\end{enumerate}

If any constraint fails:

\[
\frac{ \sigma(s_3) \neq v_3 }
{ \mathbb{C}, \sigma \Downarrow_{\text{constraints}} \text{false} }
\]

then verification fails.


\subsection{Final Rule: Complete Verification}

\[
\frac{ P \Downarrow_{\text{commit}} \text{true}, \quad (C, P) \Downarrow_{\text{map}} \sigma, \quad \mathbb{C}, \sigma \Downarrow_{\text{constraints}} \text{true} }
{ C, I, P \Downarrow \text{true} }
\]

If all three steps succeed (commitment check, value mapping, and constraint verification), then the proof is valid, and \texttt{verify!} returns true.

Otherwise:

\[
\frac{ P \Downarrow_{\text{commit}} \text{false} \quad \text{or} \quad \mathbb{C}, \sigma \Downarrow_{\text{constraints}} \text{false} }
{ C, I, P \Downarrow \text{false} }
\]

If either:

\begin{itemize}
    \item A commitment check fails, or
    \item A constraint evaluation fails,
\end{itemize}

then verification fails, and \texttt{verify!} returns false.


\end{document}
